\section*{अध्याय \ref{ch:b3:crystalline-solids} --- क्रिस्टलीय ठोस}

\begin{solution}{exo:b3:crystalline-solids:1}
(क) SC: $8\times\tfrac18 = 1$; BCC: $8\times\tfrac18 + 1 = 2$;
FCC: $8\times\tfrac18 + 6\times\tfrac12 = 4$। (ख) विकर्ण
$a\sqrt3 = 4r$: $r = a\sqrt3/4$। (ग) फलक-विकर्ण $a\sqrt2 = 4r$:
$r = a\sqrt2/4$। (घ) $0.68$ और $0.74$: दिशा-अंधा धात्विक
बंधन अधिकतम संकुलन चाहता है, और इसीलिए ताँबे, ऐलुमिनियम, चाँदी तथा
सोने की FCC (या उतनी ही सघन षट्कोणीय) संरचनाएँ।
\end{solution}

\begin{solution}{exo:b3:crystalline-solids:2}
(क) $a/\sqrt2 = \qty{2.55}{\angstrom}$। (ख) 12 --- यानी
सघन-संकुलन का समन्वय। (ग) \qty{8960}{kg/m^3}
(\cref{ex:b3:crystalline-solids:density})। (घ) $n = 4/a^3 =
\qty{8.5e28}{m^{-3}}$, अतः किसी घन माइक्रोमीटर में
$8.5\times10^{10}$ परमाणु होते हैं --- और इसीलिए चिप की धातुकर्मिकी
सांख्यिकी है, बढ़ईगीरी नहीं।
\end{solution}

\begin{solution}{exo:b3:crystalline-solids:3}
(क) $\sin\theta_1 = \lambda/2d = 0.273$: $\theta_1 =
15.8^\circ$। (ख) $n \le 2d/\lambda = 3.66$: तीन कोटियाँ। (ग)
$\sin\theta \le 1$ के कारण $n\lambda \le 2d$ होना ही चाहिए; और दृश्य
प्रकाश की $\lambda \sim \qty{5000}{\angstrom}$ हज़ार गुना लंबी है
--- कोई क्रिस्टल-तल अंतराल उसे विवर्तित नहीं कर सकता। (घ) $\sin\theta$
1\,\% गिरता है: $\theta_1 \to 15.7^\circ$। कोई विवर्तनमापी एक बढ़िया
तापमापी है --- और यही खिसकाव वह तरीक़ा है जिससे तापीय प्रसार
परमाणविक पैमाने पर नापा जाता है।
\end{solution}

\begin{solution}{exo:b3:crystalline-solids:4}
(क) $d_{100} = a > d_{110} = a/\sqrt2 > d_{111} = a/\sqrt3$। (ख)
सबसे बड़ा $d$, सबसे छोटा कोण: (100)। (ग) अष्टफलकीय (111)
कुल --- हीरे की संरचना में उसकी चादरें जोड़ों में मिलकर प्रबल रूप से
बँधी दोहरी परतें बनाती हैं, जिनके बीच चौड़े, विरल रूप से बँधे अंतराल
होते हैं: वही चिराई का तल। (घ) हर कण उसी $\theta$ पर परावर्तित करता है
पर किसी यादृच्छिक दिगंश में: परावर्तित किरणें
$2\theta$ अर्ध-कोण के शंकुओं में फैल जाती हैं, जो संसूचक को वलयों में
काटते हैं।
\end{solution}

\begin{solution}{exo:b3:crystalline-solids:5}
(क) व्यतिकरण को $\lambda$ की कोटि के पथ-अंतर चाहिए; और अंतराल
आंगस्ट्रॉम के हैं, अतः $\lambda$ को भी वही होना चाहिए। (ख) $E =
hc/\lambda \approx \qty{8}{keV}$। (ग) $p = h/\lambda$, $V =
p^2/2m_{\text{e}}e \approx \qty{60}{kV}$ (आपेक्षिकता कुछ प्रतिशत
छाँट देती है) --- यानी किसी इलेक्ट्रॉन सूक्ष्मदर्शी की कार्यकारी
वोल्टता। (घ)
$\lambda = h/\sqrt{3m_{\text{n}}k_{\text{B}}T} \approx
\qty{1.5}{\angstrom}$: कमरे के ताप के न्यूट्रॉन जन्म से ही
विवर्तन-योग्य हैं। वे नाभिकों से प्रकीर्ण होते हैं, इलेक्ट्रॉन-मेघों
से नहीं --- अतः वे हाइड्रोजन साफ़ देखते हैं और एक चुंबकीय आघूर्ण भी
ढोते हैं जो चुंबकीय क्रम मानचित्रित कर देता है; और दोनों ही
एक्स-किरणों के लिए लगभग अदृश्य हैं।
\end{solution}

\begin{solution}{exo:b3:crystalline-solids:6}
(क) आयतन $a^3$ की प्रति कोष्ठिका चार युग्म। (ख) $N_{\text{A}} =
4M/\rho a^3 = 4\times0.05844/(2165\times\qty{1.794e-28}{})
\approx \qty{6.02e23}{mol^{-1}}$। (ग) $N_{\text{A}} \propto
a^{-3}$: $a$ में 0.1\,\% की त्रुटि $N_{\text{A}}$ में 0.3\,\% है। (घ)
यह विधि परमाणु इस मान्यता पर गिनती है कि हर कोष्ठिका भरी और एक-सी है:
रिक्तियाँ, अशुद्धियाँ और पच्चीकारी-सीमाएँ, सब गिनती बिगाड़ देती हैं
--- और इसीलिए किलोग्राम की पुनर्परिभाषा के वे कट्टर हद तक निर्दोष
सिलिकॉन गोले।
\end{solution}

\begin{solution}{exo:b3:crystalline-solids:7}
(क) हर आयन $2\sum_{n\ge1}(-1)^{n+1}/n = 2\ln2$ देखता है, वह भी
$e^2/4\pi\varepsilon_0r_0$ की इकाइयों में (गुणक 2: दोनों ओर), और
वह आकर्षी है। (ख)
$e^2/4\pi\varepsilon_0r_0 = \qty{5.11}{eV}$; $\times1.7476 =
\qty{8.93}{eV}$। (ग) न्यूनतमीकरण आकर्षण का $1/9$ खा जाता है:
$U = 8.93\times\tfrac89 = \qty{7.94}{eV}$। (घ) \qty{7.9}{eV} के एक
प्रतिशत भीतर: बिंदु-आवेश \omterm{def:b3:crystalline-solids:lattice}{जालक} और एक ही कड़ाई-चरघातांक
किसी आयनी ठोस का पूरा बजट समझा देते हैं --- और क्वांटम
यांत्रिकी $r_0$ तथा उस चरघातांक के भीतर छिपी रहती है।
\end{solution}

\begin{solution}{exo:b3:crystalline-solids:8}
(क) $r_0 = 2^{1/6}\times3.40 = \qty{3.82}{\angstrom}$, जो मापे गए
\qty{3.76}{\angstrom} से 1.5\,\% ऊपर है (हर परमाणु अपने बारह
पड़ोसियों को भी अनुभव करता है, जिससे कूप कस जाता है)। (ख)
$T_{\text{m}} \sim \epsilon/2k_{\text{B}} \approx \qty{60}{K}$:
\qty{84}{K} के लिए सही पैमाना। (ग) हीलियम की \omterm{thm:b3:harmonic-oscillator:spectrum}{शून्य-बिंदु ऊर्जा}
$\hbar\omega/2$ इतने उथले कूप में स्वयं उसी कूप की टक्कर में आ जाती
है: \omterm{def:b3:crystalline-solids:lattice}{क्रिस्टल} खुद को हिलाकर बिखेर देता है, और हीलियम $T = 0$ पर भी
द्रव रहता है, जब तक उसे दबाया न जाए। (घ) प्रति बंध गहराई
$\sim\qty{0.01}{eV}$ बनाम सहसंयोजी
$\sim\qty{3.6}{eV}$: ऊर्जा में तीन सौ गुना ही पाले से हीरे तक की
पूरी दूरी है।
\end{solution}

\begin{solution}{exo:b3:crystalline-solids:9}
(क) $ka \ll 1$: $\omega \approx 2\sqrt{K/m}\,(ka/2) =
a\sqrt{K/m}\,k$। (ख) किसी कोष्ठिका को $\delta$ जितना तानिए: प्रतिबल
$\sim
K\delta/a^2$, विकृति $\delta/a$, अतः $E \sim K/a$, अर्थात् $K \sim
Ea$। (ग) $K \approx \qty{31}{N/m}$, $v = a\sqrt{K/m} \approx
\qty{4400}{m/s}$ --- मापे गए \qty{4000}{m/s} से 10\,\% भीतर, और वह
भी यंग गुणांक से अटकल लगाई किसी कमानी से। (घ) $v = a\sqrt{K/m}$:
कड़े बंध ऊपर, हलके परमाणु ऊपर --- हीरा दोनों को अधिकतम करता है, और
इसीलिए \qty{18000}{m/s} की ध्वनि तथा $\Theta_{\text{D}} \approx
\qty{2200}{K}$।
\end{solution}

\begin{solution}{exo:b3:crystalline-solids:10}
(क) $\dd\omega/\dd k \propto \cos(ka/2) = 0$, $k = \pi/a$ पर। (ख)
$u_n \propto (-1)^n$: हर परमाणु अपने दोनों पड़ोसियों से विपरीत कला
में --- यानी कोई अप्रगामी तरंग, जिसमें ऊर्जा वहीं छलकती रहती है। (ग)
$\lambda = 2a$ के साथ, क्रमिक परमाणुओं से पीछे को प्रकीर्ण तरंगों के
पथ में ठीक एक तरंगदैर्घ्य का अंतर होता है: यानी स्वयं उसी शृंखला पर
ब्रैग की शर्त। \omterm{def:b3:crystalline-solids:lattice}{जालक} अपने ही कंपन को परावर्तित कर देता है,
आगे और पीछे जाती तरंगें उस अप्रगामी प्रतिरूप में जकड़ जाती हैं, और
चलती तरंग आगे नहीं बढ़ पाती। (घ) $\omega_{\text{अधि}} = 2\sqrt{K/m} \approx
\qty{3.4e13}{rad/s}$, अर्थात् $\sim\qty{5}{THz}$: यानी दूर अवरक्त
--- \omterm{def:b3:crystalline-solids:lattice}{जालक} के कंपन और अवरक्त प्रकाश एक ही सप्तक में मिलते हैं, और यही
तथ्य अगले अभ्यास के पीछे है।
\end{solution}

\begin{solution}{exo:b3:crystalline-solids:11}
(क) प्रकाशिक मोड में $+$ और $-$ \omterm{def:b3:crystalline-solids:lattice}{उप-जालक} उलटी दिशाओं में चलते हैं:
यानी कोई दोलनशील विद्युत् द्विध्रुव, और प्रकाश-तरंग ठीक उसी को पकड़ती
है। (ख) $\omega \approx \sqrt{2K/\mu}$, जहाँ $\mu =
\qty{2.3e-26}{kg}$:
$\omega \approx \qty{4.7e13}{rad/s}$,
$\lambda = 2\pi c/\omega \approx \qty{40}{\micro m}$ --- यानी दूर
अवरक्त (नमक का मापा गया रेस्टश्ट्राहलन बैंड
$\qty{61}{\micro m}$ पर बैठता है: दो-कमानी प्रतिरूप से सही सप्तक)। (ग)
उस बैंड पर \omterm{def:b3:crystalline-solids:lattice}{क्रिस्टल} की अपनी अनुनाद तरंग को सोख लेती और फिर
परावर्तित कर देती है: पारदर्शी नमक दर्पण जैसा अपारदर्शी हो जाता है।
(घ) प्रति कोष्ठिका एक परमाणु का अर्थ है कि टकराने को कोई दूसरा
\omterm{def:b3:crystalline-solids:lattice}{उप-जालक} ही नहीं: एक ही शाखा, कोई अंतराल नहीं।
\end{solution}

\begin{solution}{exo:b3:crystalline-solids:12}
(क) $N$ परमाणु $\times$ 3 विस्थापन-दिशाएँ = $3N$
दोलक, अतः रैखिक वर्णक्रम तभी काट दिया जाता है जब वह
$3N$ मोड गिन चुके: और वही $k_{\text{D}}$ परिभाषित करता है। (ख)
$k_{\text{D}} = (6\pi^2n)^{1/3} = \qty{1.7e10}{m^{-1}}$,
$\Theta_{\text{D}} = \hbar vk_{\text{D}}/k_{\text{B}} \approx
\qty{350}{K}$ --- यानी कैलोरीमितीय \qty{343}{K}, वह भी एक ध्वनि-चाल
और एक घनत्व से। (ग) वह सीधी रेखा क्षेत्र-किनारे का चपटापन चूक जाती
है: डिबाई ऊँची आवृत्तियाँ अधिक गिन लेता है, अतः मध्यवर्ती तापों पर
अनुरूपण तनने लगता है; जबकि $T^3$ नियम (केवल लंबी तरंगें) सुरक्षित है।
(घ) दोनों यही कहते हैं कि वर्णक्रम तब समाप्त होता है जब तरंगदैर्घ्य
परमाणविक अंतराल तक पहुँच जाए --- प्रति परमाणु एक मोड, जो कभी कटाव
तरंग-सदिश के वेश में आता है और कभी कटाव-ताप के।
\end{solution}

\begin{solution}{pb:b3:crystalline-solids:1}
\textbf{1.} $2d\sin\theta = n\lambda$। एकवर्णी: एक ही
$\lambda$ हर वलय-कोण को एक ही अंतराल पर मानचित्रित कर देती है। चूर्ण:
यादृच्छिक कणों में से कुछ सदा परावर्तन के लिए सधे रहते हैं --- हर
कुल एक साथ बोल उठता है।
\textbf{2.} यहाँ $\theta = 19.05^\circ, 22.15^\circ, 32.2^\circ,
38.75^\circ$; और $\sin^2\theta = 0.107, 0.142, 0.284, 0.392$।
\textbf{3.} $d = a/\sqrt{h^2+k^2+l^2}$ को ब्रैग में रखिए
($n = 1$): $\sin^2\theta = (\lambda^2/4a^2)(h^2+k^2+l^2)$।
\textbf{4.} अनुपात $1 : 1.33 : 2.67 : 3.68$; और 3 गुना करने पर:
$3.0 : 4.0 : 8.0 : 11.0$।
\textbf{5.} सब-विषम $(111)$ तथा सब-सम $(200)$, $(220)$,
$(311)$ ठीक 3, 4, 8, 11 देते हैं; और अनुपस्थित 1, 2, 5, 6, 7 सरल
घनीय को बाहर कर देते हैं, जिसका प्रतिरूप उन्हें दिखाता।
\textbf{6.} BCC के अनुमत सम योग अनुपात
$1:2:3:4$ देते हैं --- यानी एक अलग ताल। और अनुपात इस सबके बावजूद
बचे रहते हैं कि $a$ पता न हो, तरंगदैर्घ्य खिसक जाए, या नमूना ठीक न
सधा हो: सम-विषम ज्यामिति है, अंशांकन नहीं।
\textbf{7.} थोड़े-से तल एक भोथरा व्यतिकरण-अधिकतम बनाते हैं, जैसे
पाँच झिर्रियों वाला कोई ग्रेटिंग: वलय की चौड़ाई $\sim$ क्रिस्टलिका के
आकार का व्युत्क्रम (शेरर संबंध) --- अतः तीखे वलय अच्छे उगे कणों का
प्रमाण हैं।
\textbf{8.} $a = \lambda\sqrt{h^2+k^2+l^2}/2\sin\theta = 4.088,
4.086, 4.089, 4.082\,\unit{\angstrom}$।
\textbf{9.} $a \approx \qty{4.09}{\angstrom}$।
\textbf{10.} $\rho = 4M/N_{\text{A}}a^3$।
\textbf{11.} $4/N_{\text{A}}a^3 = \qty{9.74e4}{mol/m^3}$:
चाँदी $\to \qty{10.5}{g/cm^3}$ (चाँदी के पुस्तिका-मान से मेल),
सोना $\to 19.2$ (सोने से मेल!), प्लैटिनम $\to 19.0$
(प्लैटिनम के 21.4 का खंडन --- उसका सच्चा $a$
\qty{3.92}{\angstrom} है)। प्लैटिनम बाहर; और चाँदी तथा सोना, दोनों
बचे रहते हैं, क्योंकि उनके \omterm{def:b3:crystalline-solids:lattice}{जालक नियतांक} लगभग ठीक-ठीक मिलते हैं।
\textbf{12.} द्रव्यमान पर आधारित कोई मापन: घनत्व (या बस मोलर
द्रव्यमान)। चाँदी और सोना $a$ में 0.2\,\% तक साझा करते हैं, पर उनके
परमाणविक द्रव्यमान 1.8 गुने से भिन्न हैं --- और \omterm{def:b3:crystalline-solids:lattice}{जालक} का कोई संयोग उस
खाई को नहीं पाट सकता।
\textbf{13.} ईंट पर आर्किमिडीज़: \qty{10.5}{g/cm^3} के पास
$\to$ चाँदी, जो उस जंग से भी संगत है (चाँदी काली पड़ती है; सोना
चमकता हुआ ही निकलता)। सूक्ष्म कोष्ठिका-द्रव्यमान और स्थूल तुलाई
सहमत हैं --- वही
$\rho = 4M/N_{\text{A}}a^3$, दोनों दिशाओं में पढ़ा हुआ।
\textbf{14.} परमाणु अपनी \emph{अचल औसत स्थितियों} के आसपास
खड़खड़ाते हैं: वह माध्य \omterm{def:b3:crystalline-solids:lattice}{जालक}, जो वलय-कोण तय करता है, अक्षुण्ण है;
और यादृच्छिक विस्थापन केवल तीव्रता का बँटवारा बदलते हैं।
\textbf{15.} तीव्रता $T$ के साथ चिकनाई से गिरती है (एक
$\exp(-q^2\langle x^2\rangle)$ जैसा डिबाई--वालर गुणक), और खोया
प्रकाश वलयों के बीच विसरित पृष्ठभूमि बनकर फिर प्रकट हो जाता है।
\textbf{16.} $K \approx Ea = \qty{8.3e10}{}\times
\qty{2.89e-10}{} \approx \qty{24}{N/m}$।
\textbf{17.} $v = a\sqrt{K/m} \approx \qty{3300}{m/s}$ ---
मापे गए \qty{2700}{m/s} से 20\,\% भीतर।
\textbf{18.} $\langle x^2\rangle = 3k_{\text{B}}T/K$:
वर्ग-माध्य-मूल $\approx\qty{0.23}{\angstrom}$, यानी कमरे के ताप पर
ही बंध-लंबाई का लगभग 8\,\%।
\textbf{19.} 10\,\% वाला निशान $\sim\qty{500}{K}$ तक ले जाता है,
जबकि असली \qty{1235}{K} है: सही कोटि, पर दो-सा गुणक चूक
--- हमारा एक-कमानी वाला $K$ बहुत नरम है, और लिंडेमान कोई नियम नहीं,
बस एक मापन-नियम है। सीख फिर भी टिकी रहती है: गलन तभी है जब तापीय
खड़खड़ाहट स्वयं \omterm{def:b3:crystalline-solids:lattice}{जालक} की टक्कर में आ जाए।
\textbf{20.} $d_{110} = a/\sqrt2 = \qty{2.04}{\angstrom}$:
$\sin\theta = 0.378$, $2\theta \approx 44.5^\circ$।
\textbf{21.} युग्मन की प्रबलता: इलेक्ट्रॉन (आवेशित) नैनोमीटरों में
प्रकीर्ण हो जाते हैं --- अतः सतहें और पन्नियाँ; एक्स-किरणें
माइक्रोमीटरों में --- अतः पिंडीय चूर्ण; और न्यूट्रॉन सेंटीमीटरों में
--- अतः पूरे इंजन-पुर्ज़े।
\textbf{22.} वे वलय ढहकर एक-दो चौड़े प्रभामंडल बन जाते हैं:
द्रव लघु-परास क्रम (पसंदीदा पड़ोसी-दूरी) बचाए रखता है, पर
दीर्घ-परास पंजी कोई नहीं --- तीखे व्यतिकरण के लिए कुछ आवर्ती बचा ही
नहीं।
\textbf{23.} DNA --- फ़ोटो 51। विवर्तन उन पुनरावृत्ति-दूरियों को
पढ़ लेता है जो किसी भी प्रकाश-सूक्ष्मदर्शी की पहुँच से कहीं नीचे
हैं; उस X प्रतिरूप ने कुंडली की चुगली की, और
\qty{3.4}{\angstrom} की धुँधली रेखा ने उसकी चिनी हुई सीढ़ियों की।
\textbf{24.} कि वह संरचना नियतिवादी ढंग से क्रमित है
(अर्ध-आवर्ती --- क्रमित, पर दोहराए बिना): तीखे बिंदुओं को
दीर्घ-परास कला-संसक्ति चाहिए, आवर्तिता नहीं --- और यही वह खोज है
जिसने स्वयं क्रिस्टलिकी की क्रिस्टल-परिभाषा चौड़ी कर दी।
\textbf{25.} अनुपात $3:4:8:11 \to$ FCC; $a = \qty{4.09}{\angstrom}$
$+$ घनत्व $\to$ चाँदी, न सोना न प्लैटिनम; और चूँकि गरम करना वलयों को
फीका करता है और पिघलाना उन्हें मिटा देता है, इसलिए ईंट को ठंडा ही
एक्स-किरणित कीजिए --- संग्रहालय की चाँदी, व्यतिकरण से प्रमाणित।
\end{solution}
